% Raport z benchmarku asystenta dokumentacji wizyty z lekarzem w pętli (wersja 4).
% Liczby i tabele pochodzą z fragmentów w generated/, budowanych z wyników przebiegu:
%   python benchmark.py --summary-only --export results/v4_doctor
%   python benchmark_report.py results/v4_doctor docs/benchmark/generated
% W Overleaf: wgrać ten plik razem z katalogiem generated/ (kompilator pdfLaTeX).
\documentclass[11pt,a4paper]{article}

\usepackage[utf8]{inputenc}
\usepackage[T1]{fontenc}
\usepackage[polish]{babel}
\usepackage{lmodern}
\usepackage[a4paper,margin=2.3cm]{geometry}
\usepackage{microtype}
\usepackage{booktabs}
\usepackage{longtable}
\usepackage{array}
\usepackage{enumitem}
\usepackage{xcolor}
\usepackage{pgfplots}
\pgfplotsset{compat=1.18}
\usepackage[hidelinks]{hyperref}

% generated by benchmark_report.py - do not edit
\newcommand{\nVisits}{20}
\newcommand{\nDrafts}{40}
\newcommand{\nDraftsUnchanged}{21}
\newcommand{\nCorrections}{28}
\newcommand{\nFieldsChanged}{40}
\newcommand{\nCharsEdited}{749}
\newcommand{\nDocChars}{25197}
\newcommand{\pCharsEdited}{3,0}
\newcommand{\nProductChoices}{10}
\newcommand{\nSevSignificant}{2}
\newcommand{\nSevModerate}{20}
\newcommand{\nSevCosmetic}{6}
\newcommand{\nSrcSystem}{18}
\newcommand{\nSrcDictation}{6}
\newcommand{\nSrcRegistry}{4}
\newcommand{\nDetSystemNote}{13}
\newcommand{\nDetVerifier}{1}
\newcommand{\nDetDoctor}{14}
\newcommand{\nVisitsSystemModerate}{8}
\newcommand{\nIssuesStageOne}{17}
\newcommand{\nIssuesStageOneActed}{15}
\newcommand{\nIssuesStageTwo}{3}
\newcommand{\nIssuesStageTwoActed}{1}
\newcommand{\nStageTwoSystem}{7}
\newcommand{\nStageTwoSystemByVerifier}{1}
\newcommand{\nHypotheses}{64}
\newcommand{\nHypothesesConfirmed}{57}
\newcommand{\runModel}{anthropic:claude-haiku-5-5}
\newcommand{\runRegistry}{Rejestr\_Produktow\_Leczniczych\_calosciowy\_stan\_na\_dzien\_20261110.csv}
\newcommand{\runDate}{2026-11-20}
\newcommand{\runFingerprint}{8b5b141f7eedec65}

% \input cannot be followed by \bottomrule inside a table; the TeX primitive can
\makeatletter
\newcommand*{\tabinput}[1]{\@@input #1 }
\makeatother
% \wizyta{106}: the visit number, linked to its description in the appendix
\newcommand{\wizyta}[1]{\hyperref[wizyta:#1]{#1}}

\newcolumntype{L}[1]{>{\raggedright\arraybackslash}p{#1}}
\setlength{\parskip}{0.4em}
\setlength{\parindent}{0pt}

\title{Asystent dokumentacji wizyty z~lekarzem w~pętli\\[0.3em]
       \large Benchmark interwencji lekarza -- wersja 4}
\author{}
\date{Przebieg: \runDate}

\begin{document}
\maketitle

\begin{abstract}
System na podstawie dyktowania lekarza przygotowuje projekty recepty, zaleceń dla pacjenta i~wpisu do karty wizyty;
lekarz każdy projekt poprawia i~zatwierdza, a~decyzje pozostają po jego stronie. Benchmark mierzy, \emph{jak bardzo lekarz musi ingerować}.
Na \nVisits{} scenariuszach wizyt w~podstawowej opiece zdrowotnej na trzech poziomach trudności recenzent pełniący rolę lekarza
zatwierdził bez zmian \nDraftsUnchanged{} z~\nDrafts{} projektów i~wprowadził \nCorrections{} poprawek
(\nSevSignificant{} istotne, \nSevModerate{} umiarkowanych, \nSevCosmetic{} kosmetycznych), zmieniając \pCharsEdited\,\% tekstu.
\nSrcSystem{} poprawek wynikało z~błędów systemu, \nSrcDictation{} z~niepełnego dyktowania, \nSrcRegistry{} z~ograniczeń danych rejestru leków.
Uwagi deterministyczne do recepty wskazały wszystkie braki danych w~dyktowaniu (ilość, moc), natomiast weryfikator karty wizyty wykrył
\nStageTwoSystemByVerifier{} z~\nStageTwoSystem{} błędów systemu w~karcie. Potwierdziło się \nHypothesesConfirmed{} z~\nHypotheses{}
hipotez zapisanych przed przebiegiem; najważniejsze ustalenia to brak kontroli alergii, fałszywe konflikty nazw handlowych
i~gubienie miejsca zgłoszenia przy objawach alarmowych.
\end{abstract}

\tableofcontents

% ---------------------------------------------------------------------------
\section{Cel i pytania}

Asystent ma odciążyć lekarza w~pisaniu dokumentacji, nie przejmując jego decyzji. Miarą jego użyteczności jest więc
praca, którą lekarz musi włożyć w~poprawienie projektów, oraz to, czy błędy, które przeoczyłby lekarz, mogłyby zaszkodzić pacjentowi.
Benchmark odpowiada na pytania:
\begin{enumerate}[label=P\arabic*.]
  \item Jaka część projektów nadaje się do zatwierdzenia bez zmian?
  \item Ile i~jakich poprawek wymaga lekarz (kategoria, waga) i~co spowodowało ich potrzebę: system, dyktowanie czy rejestr leków?
  \item Na ile uwagi systemu (deterministyczne do recepty i~weryfikatora do karty) wskazują lekarzowi problemy?
  \item Które cechy wizyty wiążą się z~większą liczbą poprawek?
\end{enumerate}
Benchmark nie mierzy czasu pracy ani nie porównuje wersji systemu; każdy scenariusz przebiega raz, z~lekarzem w~pętli.

% ---------------------------------------------------------------------------
\section{System}

Przebieg wizyty jest grafem LangGraph z~dwoma punktami przeglądu przez lekarza:
\begin{enumerate}[nosep]
  \item wczytanie danych pacjenta (choroby przewlekłe, leki stałe, alergie);
  \item \textbf{etap 1}: z~dyktowania równolegle agent recepty i~agent zaleceń; recepta jest dopasowywana do
        Rejestru Produktów Leczniczych, a~braki i~niejednoznaczności trafiają do deterministycznych \emph{uwag systemu};
  \item przegląd etapu 1 przez lekarza;
  \item \textbf{etap 2}: dyktowanie wywiadu i~badania, agent karty wizyty (wywiad, badanie przedmiotowe, rozpoznanie),
        weryfikator (model językowy) porównujący kartę z~dyktowaniem i~danymi pacjenta;
  \item przegląd etapu 2; lekarz wpisuje kod ICD-10 (system go nie proponuje);
  \item złożenie dokumentacji.
\end{enumerate}
Model: \texttt{\runModel}, odpowiedzi strukturalne według schematu JSON. Rejestr leków: \texttt{\runRegistry}.

\paragraph{Dopasowanie leków.} Model wyodrębnia z~dyktowania substancję, nazwę handlową, moc, postać, dawkowanie i~ilość.
Dopasowanie do rejestru jest deterministyczne: klucz pisowni łacińsko-polskiej (np.\ \emph{budesonidum}~/ \emph{budezonid}),
jednostki zapisane słownie, stężenia (\% i~mg/g), moc na dawkę inhalacyjną, grupy postaci. Gdy produkt nie jest jednoznaczny
(brak mocy, kilka postaci, sprzeczność nazwy handlowej z~substancją), system \emph{nie wybiera} i~prosi lekarza o~wybór.

\paragraph{Przegląd.} Lekarz poprawia projekt zapisany jako edytowalny JSON i~opisuje każdą poprawkę w~polu \texttt{doctor\_notes}
(kategoria, waga, źródło, uzasadnienie). Przy zatwierdzeniu sprawdzana jest poprawność dokumentu (produkt z~numerem w~rejestrze,
dawkowanie, ilość, format kodu ICD-10). Po zatwierdzeniu wizyta przechodzi do kolejnego etapu -- poprawki nie wracają do modelu.
Z~każdego etapu zostają oba pliki: projekt systemu (\texttt{stageN\_draft.json}) i~wersja zatwierdzona (\texttt{stageN\_approved.json}),
a~\texttt{log.json} zawiera listę zmienionych pól, uwagi systemu i~uwagi lekarza.

% ---------------------------------------------------------------------------
\section{Metodyka}

\subsection{Scenariusze}
\nVisits{} niezależnych wizyt w~POZ (tab.~\ref{tab:scenariusze}), każda z~dyktowaniem etapu 1 (recepta, zalecenia) i~etapu 2 (karta).
Dyktowania naśladują mowę lekarza: kolejność niezgodna z~polami dokumentów, w~poziomie C wtrącenia, autokorekty
i~błędy rozpoznawania mowy. Poziomy trudności:
\begin{itemize}[nosep]
  \item \textbf{A -- podstawowe (5)}: uporządkowane dyktowanie, 0--1 lek, do 2 cech utrudniających;
  \item \textbf{B -- typowe (9)}: naturalna kolejność, 1--5 leków (przedłużenia, preparaty złożone, inhalatory, żele, dziecko, ciąża);
  \item \textbf{C -- trudne (6)}: dyktowanie jak z~rozpoznawania mowy, wielolekowość ze zmianami leczenia, sprzeczności z~kartą pacjenta.
\end{itemize}
Każdy scenariusz ma listę \emph{cech} z~katalogu (np.\ \texttt{autokorekta}, \texttt{lateralizacja}, \texttt{konflikt-alergii},
\texttt{kontynuacja-bez-recepty}), co pozwala opisać, z~czym wiążą się poprawki.

\begin{longtable}{@{}l c L{5.6cm} L{2.9cm} L{4.6cm}@{}}
\caption{Scenariusze benchmarku.}\label{tab:scenariusze}\\
\toprule
Nr & Poz. & Przypadek & Dziedzina & Cechy \\
\midrule
\endfirsthead
\toprule
Nr & Poz. & Przypadek & Dziedzina & Cechy \\
\midrule
\endhead
\tabinput{generated/scenarios.tex}
\bottomrule
\end{longtable}

\subsection{Recenzent i procedura}
Rolę lekarza pełnił model Claude (Anthropic), pracujący w~trybie plikowym: dla każdego etapu widział dane pacjenta,
dyktowanie, projekt i~uwagi systemu, poprawiał projekt, opisywał każdą poprawkę i~zatwierdzał etap. Po obu etapach,
przed obejrzeniem zestawień, wypełniał ocenę wizyty (\texttt{analysis.json}). Recenzent nie miał wglądu w~zestawienia
w~trakcie przebiegu.

\subsection{Jednostka analizy i~klasyfikacja}
Jednostką jest \textbf{poprawka lekarza}: jedna uwaga w~\texttt{doctor\_notes}, opisująca \emph{jeden problem} projektu.
Poprawka może zmienić kilka pól (zamiana leku zmienia substancję, moc, postać, dawkowanie i~produkt), a~dwa różne problemy
w~tej samej pozycji to dwie poprawki. Każde zmienione pole jest przypisane do najbardziej szczegółowej uwagi, która je obejmuje.
Klasyfikacja według klucza (\texttt{docs/benchmark/codebook.md}):
\begin{itemize}[nosep]
  \item \textbf{kategoria}: pominięcie, treść spoza dyktowania, błędny fakt, recepta, zalecenia, forma/język;
  \item \textbf{waga} (skutek, gdyby lekarz nie poprawił): \emph{istotny} -- mógłby zaszkodzić albo zmienia kliniczny sens;
        \emph{umiarkowany} -- wymaga poprawy, szkoda mało prawdopodobna; \emph{kosmetyczny} -- bez wpływu na treść;
  \item \textbf{źródło}: \emph{system} (błąd asystenta), \emph{dyktowanie} (niepełne lub błędne dyktowanie),
        \emph{rejestr} (dane rejestru nie pozwalają wybrać produktu);
  \item \textbf{wykrycie} (wyznaczane automatycznie): \emph{uwaga systemu} lub \emph{weryfikator}, jeśli uwaga systemu
        dotyczyła któregoś z~pól poprawki, w~przeciwnym razie \emph{lekarz}.
\end{itemize}
Wpisanie kodu ICD-10 nie jest poprawką. Liczba zmienionych pól i~znaków jest podawana dodatkowo.

\subsection{Ocena wizyty}
Po każdej wizycie recenzent opisywał mocne strony i~braki projektów, wydawał werdykt
(\emph{gotowy bez zmian} / \emph{drobne poprawki} -- tylko kosmetyczne / \emph{istotne poprawki} -- co najmniej jedna
umiarkowana lub istotna / \emph{nie do użycia}) i~oceniał projekty w~skali 1--5 według pięciu kryteriów wzorowanych na
PDQI-9 (Physician Documentation Quality Instrument): dokładność, kompletność, zwięzłość, zrozumiałość dla pacjenta, spójność wewnętrzna.
Werdykt nie rozróżnia źródła poprawki -- brak ilości w~dyktowaniu też daje \emph{istotne poprawki}.

\subsection{Hipotezy}
Przed uruchomieniem dla każdego scenariusza zapisano 3--4 hipotezy: gdzie system powinien sobie poradzić, a~gdzie
spodziewano się trudności (np.\ fałszywy konflikt nazwy handlowej, brak ostrzeżenia o~alergii). Po przebiegu każdą hipotezę
sprawdzono ręcznie na projektach systemu i~zapisano dowód. Skrót SHA-256 plików scenariuszy z~chwili uruchomienia
(\texttt{\runFingerprint\ldots}) jest równy skrótowi po przebiegu.

\subsection{Rezygnacja z~oceny automatycznej}
Wcześniejsze wersje benchmarku porównywały projekty z~odpowiedzią wzorcową: oczekiwanymi lekami i~listą faktów wyszukiwanych
jako frazy. W~języku polskim takie dopasowanie jest zawodne: fakt obecny w~dokumencie w~innej formie fleksyjnej
(\emph{płukać} wobec \emph{płukanie}, \emph{kwasu foliowego} wobec \emph{kwas foliowy}) nie zostaje rozpoznany,
a~poprawna negacja (\emph{zwolnienie lekarskie nie jest potrzebne}) liczy się jako fraza niedozwolona.
Ocena wzorcowa mierzyła też system względem odpowiedzi autora scenariuszy, a~nie względem lekarza.
Dlatego jedyną miarą jest przegląd lekarza: jego poprawki, ocena wizyty i~ręcznie sprawdzone hipotezy.

% ---------------------------------------------------------------------------
\section{Przebieg prac}

\begin{enumerate}
  \item \textbf{Wersje 1--3.} Dziesięć scenariuszy, ocena automatyczna, później przegląd z~lekarzem w~pętli. Krytyczna ocena
        pokazała, że benchmark był za łatwy: dyktowania ułożone jak pola formularza, proste leki (tabletki, jedna substancja),
        brak typowych zadań POZ (przedłużenia, zmiany dawek, odstawienia), efekt sufitu. Brakowało też przypisania przyczyny
        poprawki (błąd systemu czy dyktowania) i~opisu trudności scenariuszy.
  \item \textbf{Poprawki dopasowania leków} przed nowymi scenariuszami: jednostki słownie i~duże liczby, moc na dawkę,
        stężenia, pisownia łacińsko-polska (\emph{s}/\emph{z}, \emph{-um}/\emph{-ow}), aliasy (\emph{kwas}~/ \emph{acidum});
        testy na liście około 40 leków typowych dla POZ, niezależnej od scenariuszy.
  \item \textbf{Benchmark v4}: 20 nowych scenariuszy na trzech poziomach, katalog cech, hipotezy, klucz klasyfikacji
        ze \emph{źródłem} poprawki, ocena wizyty, eksport do CSV.
  \item \textbf{Uporządkowanie metody}: (a)~usunięcie oceny automatycznej razem z~odpowiedziami wzorcowymi;
        (b)~zachowanie projektu systemu obok wersji zatwierdzonej -- lekarz edytuje projekt w~miejscu, więc po zatwierdzeniu
        projekt jest przywracany, a~wersja lekarza zapisywana osobno; (c)~poprawka lekarza jako jednostka analizy
        z~zasadą \emph{jedna uwaga = jeden problem} zamiast liczenia zmienionych pól (zamiana antybiotyku to jedna poprawka,
        nie sześć zmian). Usunięto też pliki i~kod, które służyły tylko ocenie automatycznej.
  \item \textbf{Przebieg benchmarku} (\runDate): 20 wizyt, 40 przeglądów, eksport danych i~ten raport.
\end{enumerate}

% ---------------------------------------------------------------------------
\section{Wyniki}

\subsection{Skala interwencji}
Bez zmian zatwierdzono \nDraftsUnchanged{} z~\nDrafts{} projektów. Lekarz wprowadził \nCorrections{} poprawek
(\nFieldsChanged{} zmienionych pól, \nCharsEdited{} z~\nDocChars{} znaków, czyli \pCharsEdited\,\%), w~tym \nProductChoices{}
wyborów produktu z~listy kandydatów. Poprawkę wynikającą z~błędu systemu o~wadze co najmniej umiarkowanej
wymagało \nVisitsSystemModerate{} z~\nVisits{} wizyt. Liczba poprawek rośnie z~poziomem trudności (tab.~\ref{tab:poziomy},
rys.~\ref{rys:poziomy}); poziom A pozostaje blisko sufitu i~pełni rolę kontrolną.

\begin{table}[ht]
\centering
\caption{Interwencje według poziomu trudności.}\label{tab:poziomy}
\small
\begin{tabular}{@{}l r r r r r r r r@{}}
\toprule
Poziom & Wizyty & Bez zmian & Poprawki & na wizytę & z~winy systemu & na wizytę & Istotne & Pola \\
\midrule
\tabinput{generated/levels.tex}
\bottomrule
\end{tabular}
\end{table}

\begin{figure}[ht]
\centering
\begin{tikzpicture}
\begin{axis}[ybar stacked, bar width=22pt, width=9cm, height=5.5cm,
             symbolic x coords={A,B,C}, xtick=data, ymin=0,
             ylabel={poprawki na wizytę}, xlabel={poziom trudności},
             legend style={at={(1.02,0.5)}, anchor=west, draw=none}, legend cell align=left]
\tabinput{generated/level_chart.tex}
\end{axis}
\end{tikzpicture}
\caption{Średnia liczba poprawek na wizytę według poziomu i~źródła.}\label{rys:poziomy}
\end{figure}

\subsection{Rodzaj i~waga poprawek}
Większość poprawek dotyczyła recepty (tab.~\ref{tab:kategorie}): wyboru produktu, którego system słusznie nie wybrał sam,
brakującej ilości z~dyktowania, a~w~dwóch przypadkach -- złego leku. Obie poprawki istotne dotyczyły recepty:
w~wizycie~\wizyta{106} na recepcie znalazł się lek kontynuowany bez zmian zamiast nowo dołożonego (błąd systemu),
w~wizycie~\wizyta{116} lekarz podyktował antybiotyk, na który pacjentka jest uczulona, a~system tego nie zasygnalizował.
W~zaleceniach dominowały pominięcia (miejsce zgłoszenia przy objawach alarmowych, zmiany leczenia), w~karcie -- drobne
dopiski spoza dyktowania i~pojedyncze zmiany sensu.

\begin{table}[ht]
\centering
\caption{Poprawki według kategorii i~wagi.}\label{tab:kategorie}
\begin{tabular}{@{}l r r r r@{}}
\toprule
Kategoria & istotny & umiarkowany & kosmetyczny & razem \\
\midrule
\tabinput{generated/categories.tex}
\bottomrule
\end{tabular}
\end{table}

\subsection{Źródło poprawek i~uwagi systemu}
\nSrcSystem{} poprawek wynikało z~błędów systemu, \nSrcDictation{} z~dyktowania, \nSrcRegistry{} z~danych rejestru
(tab.~\ref{tab:zrodla}). Deterministyczne uwagi do recepty działały dobrze: z~\nIssuesStageOne{} uwag \nIssuesStageOneActed{}
doprowadziło do zmiany (pozostałe dwie to uwagi o~mocy, rozwiązane razem z~wyborem produktu), a~każdy brak w~dyktowaniu
(ilość, moc) został wskazany. Słabym ogniwem jest weryfikator karty: zgłosił \nIssuesStageTwo{} uwagi, z~których
\nIssuesStageTwoActed{} była trafna, i~wykrył \nStageTwoSystemByVerifier{} z~\nStageTwoSystem{} błędów systemu w~karcie.
Błędy, których nie wskazała żadna uwaga, znalazł lekarz (\nDetDoctor{} poprawek) -- w~tym błąd istotny z~wizyty~\wizyta{116}.
Żaden produkt wybrany przez system samodzielnie nie okazał się błędny: każda niejednoznaczność skończyła się prośbą o~wybór.

\begin{table}[ht]
\centering
\caption{Poprawki według źródła i~tego, kto wskazał problem.}\label{tab:zrodla}
\begin{tabular}{@{}l r r r r@{}}
\toprule
Źródło & uwaga systemu & weryfikator & lekarz & razem \\
\midrule
\tabinput{generated/sources.tex}
\bottomrule
\end{tabular}
\end{table}

\subsection{Cechy wizyt}
Tabela~\ref{tab:cechy} zestawia średnią liczbę poprawek w~wizytach z~daną cechą. Najwięcej poprawek z~winy systemu
towarzyszyło preparatom złożonym i~nazwom handlowym (fałszywe konflikty), wiekowi podeszłemu i~wielolekowości
(zmiany leczenia w~zaleceniach), kontynuacji leku bez recepty. Cechy kojarzone z~ryzykiem błędu klinicznego --
lateralizacja, dawkowanie wagowe u~dzieci, autokorekty, zaprzeczenia -- same nie wymagały ani jednej poprawki: strony,
dawki wagowe, autokorekty i~zaprzeczenia zostały oddane bezbłędnie, a~poprawki w~tych wizytach dotyczyły innych elementów
(np.\ miejsca zgłoszenia przy objawach alarmowych).
Ze względu na małą liczbę wizyt na cechę (1--10) i~powiązanie cech z~poziomem zestawienie jest wyłącznie opisowe.

\begin{longtable}{@{}l r r r@{}}
\caption{Poprawki w~wizytach z~daną cechą (średnio na wizytę), malejąco według poprawek z~winy systemu.}\label{tab:cechy}\\
\toprule
Cecha & Wizyty & Poprawki & z~winy systemu \\
\midrule
\endfirsthead
\toprule
Cecha & Wizyty & Poprawki & z~winy systemu \\
\midrule
\endhead
\tabinput{generated/features.tex}
\bottomrule
\end{longtable}

\subsection{Ocena wizyt}
\begin{table}[ht]
\centering
\caption{Średnie oceny projektów (1--5) i~werdykty według poziomu.}\label{tab:oceny}
\small
\begin{tabular}{@{}l r r r r r r r r r@{}}
\toprule
 & \multicolumn{5}{c}{oceny 1--5} & \multicolumn{4}{c}{werdykty} \\
\cmidrule(lr){2-6}\cmidrule(l){7-10}
Poziom & dokł. & kompl. & zwięzł. & zrozum. & spójn. & bez zmian & drobne & istotne & nie do użycia \\
\midrule
\tabinput{generated/ratings.tex}
\bottomrule
\end{tabular}
\end{table}
Projekty oceniono wysoko (tab.~\ref{tab:oceny}); najniżej wypadła dokładność, obniżana przez pojedyncze zmiany sensu i~zły lek
na recepcie. Żadnej wizyty nie oceniono jako \emph{nie do użycia}. Werdykt \emph{istotne poprawki} obejmuje także wizyty,
w~których jedyną poprawką był brak ilości w~dyktowaniu lub wybór postaci, której rejestr nie rozstrzyga.

\subsection{Hipotezy}
Potwierdziło się \nHypothesesConfirmed{} z~\nHypotheses{} hipotez (tab.~\ref{tab:hipotezy}). Nie potwierdziły się:
dopasowanie mocy zapisanej słownie (101), lek kontynuowany poza receptą (106), automatyczny wybór Symbicortu i~Xonvei (108, 113),
pominięcie przedłużanych leków (107 -- przewidywane ryzyko się nie spełniło) oraz przewidywanie, że karta przemilczy sprzeczność
w~wywiadzie o~alergii (116 -- karta ją odnotowała).

\small
\begin{longtable}{@{}l L{6.1cm} c L{6.2cm}@{}}
\caption{Hipotezy zapisane przed przebiegiem i~ich weryfikacja.}\label{tab:hipotezy}\\
\toprule
Id & Hipoteza & Potw. & Dowód \\
\midrule
\endfirsthead
\toprule
Id & Hipoteza & Potw. & Dowód \\
\midrule
\endhead
\tabinput{generated/hypotheses.tex}
\bottomrule
\end{longtable}
\normalsize

% ---------------------------------------------------------------------------
\section{Ustalenia}

\begin{enumerate}[label=U\arabic*.]
  \item \textbf{Brak kontroli alergii} (wizyta~\wizyta{116}). Antybiotyk sprzeczny z~alergią w~karcie pacjentki przeszedł
        bez żadnej uwagi systemu; to jedyny błąd, który mógłby bezpośrednio zaszkodzić pacjentowi.
        Zalecenie: deterministyczna kontrola substancji z~recepty względem alergii w~danych pacjenta (z~grupami leków,
        np.\ penicyliny).
  \item \textbf{Zły lek na recepcie przy leku kontynuowanym} (\wizyta{106}). Lek \emph{zostawiamy bez zmian} trafił na receptę,
        a~nowy lek z~niej wypadł, choć zalecenia o~nim mówią. W~pozostałych wizytach z~lekiem kontynuowanym (104, 113, 117) reguła
        działała. Zalecenie: kontrola spójności recepty z~zaleceniami (lek wymieniony w~zaleceniach jako nowy, a~nieobecny na recepcie).
  \item \textbf{Fałszywe konflikty nazwy handlowej} (\wizyta{108}, \wizyta{113}). Substancje połączone spójnikiem
        \emph{z} (\emph{budezonid z~formoterolem}) nie są rozpoznawane jako skład preparatu złożonego. Błąd wystąpił w~obu takich
        wizytach i~jest łatwy do usunięcia: normalizacja \emph{X z~Y} do \emph{X + Y} przed porównaniem.
  \item \textbf{Moc zapisana słownie lub bez jednostki} (\wizyta{101}). Model niekonsekwentnie
        zamienia liczebniki na liczby; parser rejestru wymaga liczby z~jednostką. Zalecenie: liczebniki w~parserze mocy
        albo wymóg zapisu cyfrowego z~jednostką w~schemacie.
  \item \textbf{Gubienie miejsca zgłoszenia przy objawach alarmowych} (\wizyta{108}, \wizyta{109}, \wizyta{120}).
        Gdy objawy są rozbite na osobne punkty, podyktowane \emph{pilnie na SOR / do szpitala} znika. Zalecenie: osobne pole
        na miejsce i~pilność zgłoszenia w~schemacie zaleceń.
  \item \textbf{Zmiany sensu} (\wizyta{111}, \wizyta{117}): przeniesienie oceny działania leku z~pacjenta na lekarza,
        dawka 402~mg z~\emph{czterysta, dwa, trzy razy dziennie}, a~także drobne dopiski spoza dyktowania (skróty, jednostki,
        oceny typu \emph{stan ogólny dobry}).
  \item \textbf{Zalecenia przy wielolekowości} (\wizyta{117}): zmiany leczenia (czasowe zwiększenie dawki, nowy lek,
        kontynuacja leku poza receptą) nie trafiają do zaleceń dla pacjenta.
  \item \textbf{Weryfikator nie jest siatką bezpieczeństwa}: wykrył \nStageTwoSystemByVerifier{} z~\nStageTwoSystem{} błędów
        w~karcie, a~część jego uwag to fałszywe alarmy (np.\ błędne odczytanie autokorekty lekarza w~116).
  \item \textbf{Mocne strony}: bez błędów w~lateralizacji, wartościach liczbowych (dwie temperatury, dwie masy ciała, wyniki
        laboratoryjne), dawkowaniu wagowym u~dzieci (ostateczne 10,5~ml mimo liczenia na głos), zaprzeczeniach, autokorektach
        i~wtrąceniach; pilne skierowania (110, 120) oddane wprost; żaden produkt nie został wybrany błędnie bez wiedzy lekarza.
\end{enumerate}

% ---------------------------------------------------------------------------
\section{Ograniczenia}

\begin{itemize}
  \item \textbf{Brak miary zgodności.} Nie policzono zgodności klasyfikacji między recenzentami ($\kappa$ Cohena);
        klucz klasyfikacji powstał z~myślą o~takim pomiarze.
  \item \textbf{Mała próba i~jeden przebieg.} \nVisits{} wizyt, każda raz; model językowy nie jest deterministyczny,
        więc wynik pojedynczej wizyty może się różnić między przebiegami. Analiza cech jest opisowa (1--10 wizyt na cechę, cechy powiązane z~poziomem).
  \item \textbf{Warunki sztuczne.} Dyktowania są tekstem (bez rozpoznawania mowy), scenariusze i~dane pacjentów są fikcyjne,
        użyto jednego modelu językowego.
  \item \textbf{Werdykt} nie rozróżnia źródła poprawek, więc zawyża udział wizyt z~\emph{istotnymi poprawkami};
        właściwszą miarą jakości systemu jest liczba poprawek z~winy systemu (\nSrcSystem{}).
\end{itemize}

% ---------------------------------------------------------------------------
\section{Wnioski i~dalsze kroki}

Asystent przygotowuje projekty, które w~połowie przypadków można zatwierdzić bez zmian, a~pozostałe wymagają zwykle jednej--dwóch
poprawek, w~większości umiarkowanych i~wskazanych przez sam system. Projekty wiernie oddają dyktowanie, także chaotyczne,
z~autokorektami i~liczeniem dawek na głos. Lekarz pozostaje jednak niezbędny: lek sprzeczny z~alergią (\wizyta{116}) znalazł sam,
przy złym leku (\wizyta{106}) uwagi systemu wskazywały jedynie brak mocy i~ilości, a~nie sam błąd, a~weryfikator karty wykrywa niewiele. Najbardziej obiecujące zmiany systemu to deterministyczne
kontrole: alergii względem karty pacjenta, spójności recepty z~zaleceniami, preparatów złożonych zapisanych ze spójnikiem
\emph{z} oraz mocy zapisanej słownie. Kolejny etap badania powinien obejmować przegląd części wizyt przez lekarza
(zgodność klasyfikacji) i~kilka przebiegów każdego scenariusza (zmienność modelu).

% ---------------------------------------------------------------------------
\appendix
\section{Opisy wizyt}
Dla każdej wizyty: scenariusz, ocena recenzenta (podsumowanie, werdykt, oceny), mocne strony, braki i~wszystkie poprawki
z~klasyfikacją (kategoria / waga / źródło / kto wskazał) i~uzasadnieniem lekarza.

{\sloppy
% generated by benchmark_report.py - do not edit
\subsection*{101. Ostre paciorkowcowe zapalenie gardła i migdałków}\phantomsection\addcontentsline{toc}{subsection}{101. Ostre paciorkowcowe zapalenie gardła i migdałków}\label{wizyta:101}
\textbf{Poziom} A · choroby zakaźne / laryngologia · \textbf{cechy:} brak-ilosci\par
\textit{Scenariusz.} 28-latek, Centor/McIsaac 4 pkt, dodatni szybki test na antygen paciorkowca. Fenoksymetylpenicylina przez 10 dni, paracetamol doraźnie. Lekarz nie podaje ilości antybiotyku.\par
\textit{Ocena recenzenta.} Zalecenia i karta wierne dyktowaniu; recepta wymagała ręcznego wyboru antybiotyku, bo model zapisał moc słownie („1 milion jednostek”), a dyktowanie nie podało ilości. Werdykt: \textbf{istotne poprawki}; oceny 1--5: dokładność 4, kompletność 5, zwięzłość 4, zrozumiałość dla pacjenta 5, spójność wewnętrzna 5.\par
\textit{Mocne strony:}
\begin{itemize}[nosep]
\item Brak ilości fenoksymetylpenicyliny oznaczony przez system.
\item Objawy alarmowe (szczękościsk, jednostronny obrzęk gardła, ślinienie, duszność) przepisane kompletnie.
\item Karta: obie temperatury (39 w wywiadzie, 38,7 w badaniu) i skala Centora przypisane poprawnie.
\end{itemize}
\textit{Braki:}
\begin{itemize}[nosep]
\item Moc zapisana słownie – rejestr nie znalazł produktu mimo jednoznacznego dyktowania.
\item Dopisany skrót „(GAS)” w badaniu, którego nie dyktowano; weryfikator go nie zgłosił.
\end{itemize}
\textit{Poprawki lekarza:}
\begin{itemize}[nosep]
\item etap 1 (recepta i zalecenia), \texttt{prescription.0}: recepta / umiarkowany / źródło: system / wykrył: uwaga systemu. Model zapisał moc słownie („1 milion jednostek”), przez co rejestr nie znalazł produktu; poprawiam moc na 1 000 000 j.m. i wybieram Ospen 1000.
\item etap 1 (recepta i zalecenia), \texttt{prescription.0.quantity}: recepta / umiarkowany / źródło: dyktowanie / wykrył: uwaga systemu. Nie podałem ilości (system to oznaczył); 3 $\times$ 1 tabletka przez 10 dni = 30 tabletek.
\item etap 2 (karta wizyty), \texttt{visit\_note.physical\_examination}: treść spoza dyktowania / kosmetyczny / źródło: system / wykrył: lekarz. Dopisany skrót „(GAS)”, którego nie dyktowałem.
\end{itemize}

\subsection*{102. Bakteryjne zapalenie spojówek obu oczu}\phantomsection\addcontentsline{toc}{subsection}{102. Bakteryjne zapalenie spojówek obu oczu}\label{wizyta:102}
\textbf{Poziom} A · okulistyka · \textbf{cechy:} postac-nietabletkowa, zaprzeczenia\par
\textit{Scenariusz.} 37-latka nosząca soczewki kontaktowe, ropna wydzielina obu oczu bez bólu i światłowstrętu. Lewofloksacyna w kroplach do oczu ze schematem malejącym.\par
\textit{Ocena recenzenta.} Oba projekty zatwierdzone bez zmian: krople do oczu dopasowane do właściwej postaci, schemat malejący i zaprzeczenia zachowane. Werdykt: \textbf{gotowy bez zmian}; oceny 1--5: dokładność 5, kompletność 5, zwięzłość 5, zrozumiałość dla pacjenta 5, spójność wewnętrzna 5.\par
\textit{Mocne strony:}
\begin{itemize}[nosep]
\item Lewofloksacyna dopasowana do kropli do oczu 5 mg/ml, nie do innej postaci.
\item Schemat dawkowania (co 2 godziny przez 2 dni, potem 4 razy dziennie, łącznie 7 dni) przepisany dokładnie.
\item Zaprzeczenia (bez bólu oka, bez światłowstrętu) i kolejność zajęcia oczu zachowane w karcie.
\end{itemize}
\textit{Poprawki lekarza:} brak – oba projekty zatwierdzone bez zmian.\par

\subsection*{103. Ostre zapalenie oskrzeli bez antybiotyku}\phantomsection\addcontentsline{toc}{subsection}{103. Ostre zapalenie oskrzeli bez antybiotyku}\label{wizyta:103}
\textbf{Poziom} A · pulmonologia / infekcje · \textbf{cechy:} bez-lekow, zaprzeczenia\par
\textit{Scenariusz.} 41-letni palacz z ostrym, najpewniej wirusowym zapaleniem oskrzeli. Lekarz świadomie nie przepisuje antybiotyku; paracetamol pacjent ma w domu.\par
\textit{Ocena recenzenta.} Oba projekty bez zmian: pusta recepta mimo wzmianki o paracetamolu, zaprzeczenia i niepewność rozpoznania zachowane. Werdykt: \textbf{gotowy bez zmian}; oceny 1--5: dokładność 5, kompletność 5, zwięzłość 5, zrozumiałość dla pacjenta 5, spójność wewnętrzna 5.\par
\textit{Mocne strony:}
\begin{itemize}[nosep]
\item Recepta pusta – paracetamol „który pacjent ma w domu” tylko w zaleceniach.
\item Zaprzeczenia (bez duszności, bólu w klatce, krwioplucia, trzeszczeń) zachowane.
\item Rozpoznanie z zastrzeżeniem „najpewniej wirusowe”; dwie temperatury przypisane poprawnie.
\end{itemize}
\textit{Poprawki lekarza:} brak – oba projekty zatwierdzone bez zmian.\par

\subsection*{104. Nadciśnienie tętnicze niewyrównane – zwiększenie dawki amlodypiny}\phantomsection\addcontentsline{toc}{subsection}{104. Nadciśnienie tętnicze niewyrównane – zwiększenie dawki amlodypiny}\label{wizyta:104}
\textbf{Poziom} A · kardiologia · \textbf{cechy:} zmiana-dawki, kontynuacja-bez-recepty\par
\textit{Scenariusz.} 65-latka z nadciśnieniem niewyrównanym w pomiarach domowych i gabinetowych. Amlodypina 5 $\rightarrow$ 10 mg, atorwastatyna bez zmian (bez recepty), stara dawka do odstawienia.\par
\textit{Ocena recenzenta.} Oba projekty bez zmian: na recepcie tylko nowa dawka amlodypiny, zmiana dawki jasno wyjaśniona pacjentce, pomiary domowe i gabinetowe rozdzielone. Werdykt: \textbf{gotowy bez zmian}; oceny 1--5: dokładność 5, kompletność 5, zwięzłość 5, zrozumiałość dla pacjenta 5, spójność wewnętrzna 5.\par
\textit{Mocne strony:}
\begin{itemize}[nosep]
\item Atorwastatyna „bez zmian” poza receptą, ujęta w zaleceniach.
\item Odstawienie dotychczasowej amlodypiny 5 mg wprost w zaleceniach.
\item Pomiar domowy (156/94) w wywiadzie, dwa pomiary gabinetowe w badaniu.
\end{itemize}
\textit{Poprawki lekarza:} brak – oba projekty zatwierdzone bez zmian.\par

\subsection*{105. Niedoczynność tarczycy w przebiegu choroby Hashimoto}\phantomsection\addcontentsline{toc}{subsection}{105. Niedoczynność tarczycy w przebiegu choroby Hashimoto}\label{wizyta:105}
\textbf{Poziom} A · endokrynologia · \textbf{cechy:} jednostki-slownie, wartosci-liczbowe\par
\textit{Scenariusz.} 47-latka z objawami niedoczynności, TSH 12,4, FT4 poniżej normy, anty-TPO 450. Lewotyroksyna 50 mikrogramów, kontrola TSH za 6–8 tygodni.\par
\textit{Ocena recenzenta.} Oba projekty bez zmian: „pięćdziesiąt mikrogramów” dopasowane do lewotyroksyny 50 mcg, wyniki laboratoryjne i termin kontroli przepisane bez błędów. Werdykt: \textbf{gotowy bez zmian}; oceny 1--5: dokładność 5, kompletność 5, zwięzłość 5, zrozumiałość dla pacjenta 5, spójność wewnętrzna 5.\par
\textit{Mocne strony:}
\begin{itemize}[nosep]
\item Moc w mikrogramach zapisana słownie poprawnie dopasowana do produktu (Eltroxin 50 mcg).
\item TSH 12,4, FT4 0,7 i anty-TPO 450 z jednostkami, bez pomyłek.
\item Kolejność kontroli (TSH za 6–8 tygodni, potem wizyta z wynikiem) zachowana.
\end{itemize}
\textit{Poprawki lekarza:} brak – oba projekty zatwierdzone bez zmian.\par

\subsection*{106. Cukrzyca typu 2 niewyrównana – dołączenie empagliflozyny}\phantomsection\addcontentsline{toc}{subsection}{106. Cukrzyca typu 2 niewyrównana – dołączenie empagliflozyny}\label{wizyta:106}
\textbf{Poziom} B · diabetologia · \textbf{cechy:} kontynuacja-bez-recepty, wartosci-liczbowe, zaprzeczenia\par
\textit{Scenariusz.} 61-latek z cukrzycą typu 2, HbA1c 8,1\% (wcześniej 7,6\%). Metformina bez zmian, nowa empagliflozyna, zasady postępowania w dniach choroby, skierowanie na badanie dna oka. Dyktowanie w naturalnej kolejności, bez nagłówków.\par
\textit{Ocena recenzenta.} Zalecenia i karta poprawne, ale recepta zawierała zły lek: kontynuowaną bez zmian metforminę zamiast nowo dołożonej empagliflozyny, której na recepcie w ogóle nie było. Werdykt: \textbf{istotne poprawki}; oceny 1--5: dokładność 2, kompletność 3, zwięzłość 5, zrozumiałość dla pacjenta 5, spójność wewnętrzna 2.\par
\textit{Mocne strony:}
\begin{itemize}[nosep]
\item Zalecenia kompletne, z zasadą odstawienia empagliflozyny przy wymiotach i biegunce oraz skierowaniem do okulisty.
\item Wartości HbA1c (8,1\% obecnie, 7,6\% w marcu), eGFR i glikemie przepisane bez pomyłek.
\item Uwagi systemu (brak mocy i ilości metforminy) zwróciły uwagę na błędną pozycję recepty.
\end{itemize}
\textit{Braki:}
\begin{itemize}[nosep]
\item Recepta odwrócona względem dyktowania: lek „bez zmian” na recepcie, nowy lek pominięty – niespójność z zaleceniami, które mówią o recepcie na empagliflozynę.
\item Łacińska pisownia „perindopryl” w wywiadzie.
\end{itemize}
\textit{Poprawki lekarza:}
\begin{itemize}[nosep]
\item etap 1 (recepta i zalecenia), \texttt{prescription.0}: recepta / istotny / źródło: system / wykrył: uwaga systemu. Na recepcie znalazła się metformina („zostawiamy, bez zmian”), a nowo dołożona empagliflozyna 10 mg (3 $\times$ 30 tabletek) w ogóle nie trafiła na receptę, choć zalecenia o niej mówią. Zły lek na recepcie – zamieniam na empagliflozynę.
\item etap 2 (karta wizyty), \texttt{visit\_note.history}: forma/język / kosmetyczny / źródło: system / wykrył: lekarz. „perindopryl” zamiast polskiej nazwy „peryndopryl” (tak w dyktowaniu i w karcie pacjenta).
\end{itemize}

\subsection*{107. Przedłużenie leków przewlekłych u 74-latki z wielochorobowością}\phantomsection\addcontentsline{toc}{subsection}{107. Przedłużenie leków przewlekłych u 74-latki z wielochorobowością}\label{wizyta:107}
\textbf{Poziom} B · medycyna rodzinna / geriatria · \textbf{cechy:} polipragmazja, przedluzenie-recepty, odstawienie, wiek-podeszly, jednostki-slownie\par
\textit{Scenariusz.} Wizyta po recepty na stałe leki: 5 leków na 3 miesiące z różnymi ilościami opakowań, odstawienie ketoprofenu z powodu PChN. Sprawdza, czy model nie pomija leków „przedłużanych” (reguła o lekach kontynuowanych).\par
\textit{Ocena recenzenta.} Przedłużenie pięciu leków przewlekłych bez błędów modelu: wszystkie na recepcie z właściwymi ilościami, ketoprofen i paracetamol poza receptą; jedyne interwencje to wybór postaci dwóch leków, których rejestr nie rozstrzyga. Werdykt: \textbf{istotne poprawki}; oceny 1--5: dokładność 5, kompletność 5, zwięzłość 5, zrozumiałość dla pacjenta 4, spójność wewnętrzna 5.\par
\textit{Mocne strony:}
\begin{itemize}[nosep]
\item Reguła o lekach kontynuowanych nie spowodowała pominięcia przedłużanych leków (wszystkie 5 na recepcie).
\item Ilości opakowań przypisane do właściwych pozycji.
\item Ketoprofen jako zalecenie odstawienia, paracetamol tylko w zaleceniach; karta z kompletem wyników laboratoryjnych.
\end{itemize}
\textit{Braki:}
\begin{itemize}[nosep]
\item Dwie pozycje wymagały ręcznego wyboru postaci (metformina, lewotyroksyna) – ograniczenie danych rejestru przy niepodanej postaci.
\end{itemize}
\textit{Poprawki lekarza:}
\begin{itemize}[nosep]
\item etap 1 (recepta i zalecenia), \texttt{prescription.3.product}: recepta / umiarkowany / źródło: rejestr / wykrył: uwaga systemu. Metformina 1000 mg jest w rejestrze jako tabletki powlekane i tabletki o przedłużonym uwalnianiu; postaci nie podałam – wybieram tabletki powlekane (dotychczasowe leczenie 2 $\times$ dziennie).
\item etap 1 (recepta i zalecenia), \texttt{prescription.4.product}: recepta / umiarkowany / źródło: rejestr / wykrył: uwaga systemu. Lewotyroksyna 75 mcg jest w rejestrze jako tabletki i roztwór doustny; wybieram tabletki.
\end{itemize}

\subsection*{108. Zaostrzenie astmy oskrzelowej}\phantomsection\addcontentsline{toc}{subsection}{108. Zaostrzenie astmy oskrzelowej}\label{wizyta:108}
\textbf{Poziom} B · pulmonologia · \textbf{cechy:} lek-zlozony, nazwa-handlowa, postac-nietabletkowa, jednostki-slownie\par
\textit{Scenariusz.} 25-latka z astmą, umiarkowane zaostrzenie po infekcji (PEF 65\%). Symbicort w zwiększonej dawce z powrotem do dawki podstawowej, salbutamol doraźnie, prednizon przez 5 dni. Uczulenie na kwas acetylosalicylowy w karcie.\par
\textit{Ocena recenzenta.} Leki i schematy dawkowania poprawne (prednizon, nie prednizolon; schemat Symbicortu), ale Symbicort wymagał ręcznego wyboru po fałszywym konflikcie nazwy handlowej, z objawów alarmowych zniknęło „pilnie na SOR”, a z wywiadu alergia na kwas acetylosalicylowy. Werdykt: \textbf{istotne poprawki}; oceny 1--5: dokładność 4, kompletność 3, zwięzłość 5, zrozumiałość dla pacjenta 4, spójność wewnętrzna 5.\par
\textit{Mocne strony:}
\begin{itemize}[nosep]
\item Salbutamol dopasowany do aerozolu 100 mcg/dawkę, prednizon do prednizonu 20 mg.
\item Schemat „2 inhalacje 2 $\times$ dziennie przez 2 tygodnie, potem 1” zachowany.
\item Karta z kompletem parametrów zaostrzenia (PEF 280 l/min, 65\% najlepszej wartości, saturacja 96\%).
\end{itemize}
\textit{Braki:}
\begin{itemize}[nosep]
\item Fałszywy konflikt nazwy handlowej (spójnik „z” w „budezonid z formoterolem”).
\item Pominięte miejsce zgłoszenia przy objawach alarmowych (SOR).
\item Pominięta alergia na kwas acetylosalicylowy z karty pacjentki.
\end{itemize}
\textit{Poprawki lekarza:}
\begin{itemize}[nosep]
\item etap 1 (recepta i zalecenia), \texttt{prescription.0.product}: recepta / umiarkowany / źródło: system / wykrył: uwaga systemu. Fałszywy konflikt nazwy handlowej: substancje zapisane „budezonid z formoterolem” (spójnik „z”), a Symbicort to właśnie budezonid + formoterol. Wybieram 160/4,5 mcg.
\item etap 1 (recepta i zalecenia), \texttt{patient\_instructions.recommendations.1}: treść spoza dyktowania / kosmetyczny / źródło: system / wykrył: lekarz. Podyktowałem „unikać dymu”; model zawęził do „dymu papierosowego”.
\item etap 1 (recepta i zalecenia), \texttt{patient\_instructions.warning\_signs}: pominięcie / umiarkowany / źródło: system / wykrył: lekarz. Pominięte „pilnie na SOR” – pacjentka nie wie, gdzie się zgłosić przy objawach zaostrzenia.
\item etap 2 (karta wizyty), \texttt{visit\_note.history}: pominięcie / umiarkowany / źródło: system / wykrył: lekarz. Pominięta alergia na kwas acetylosalicylowy z karty pacjentki (duszność) – istotna przy zaostrzeniu astmy.
\end{itemize}

\subsection*{109. Ostre odmiedniczkowe zapalenie nerek prawostronne}\phantomsection\addcontentsline{toc}{subsection}{109. Ostre odmiedniczkowe zapalenie nerek prawostronne}\label{wizyta:109}
\textbf{Poziom} B · nefrologia / infekcje · \textbf{cechy:} wartosci-liczbowe, zaprzeczenia, lateralizacja\par
\textit{Scenariusz.} 33-latka z gorączką, dreszczami i bólem prawej okolicy lędźwiowej, dodatni Goldflam po prawej, CRP 128, test ciążowy ujemny. Cyprofloksacyna ambulatoryjnie, posiew, kontrola za 3 dni.\par
\textit{Ocena recenzenta.} Recepta i karta bez zarzutu (strona prawa spójna, test ciążowy odnotowany); w objawach alarmowych ponownie zabrakło podyktowanego „pilnie do szpitala”. Werdykt: \textbf{istotne poprawki}; oceny 1--5: dokładność 5, kompletność 4, zwięzłość 5, zrozumiałość dla pacjenta 4, spójność wewnętrzna 5.\par
\textit{Mocne strony:}
\begin{itemize}[nosep]
\item Cyprofloksacyna i paracetamol dopasowane poprawnie, z ilością 14 tabletek.
\item Strona prawa (Goldflam dodatni, ból) i lewa (Goldflam ujemny) bez zamiany; ujemny test ciążowy w wywiadzie.
\item Zalecenie o odstępie od wapnia, magnezu i żelaza zachowane.
\end{itemize}
\textit{Braki:}
\begin{itemize}[nosep]
\item Pominięte miejsce zgłoszenia przy objawach alarmowych (szpital) – ten sam błąd co przy astmie (108).
\end{itemize}
\textit{Poprawki lekarza:}
\begin{itemize}[nosep]
\item etap 1 (recepta i zalecenia), \texttt{patient\_instructions.warning\_signs}: pominięcie / umiarkowany / źródło: system / wykrył: lekarz. Pominięte „pilnie do szpitala” – przy odmiedniczkowym zapaleniu nerek pacjentka musi wiedzieć, że chodzi o szpital, nie o kontrolę.
\end{itemize}

\subsection*{110. Podejrzenie zakrzepicy żył głębokich lewej łydki}\phantomsection\addcontentsline{toc}{subsection}{110. Podejrzenie zakrzepicy żył głębokich lewej łydki}\label{wizyta:110}
\textbf{Poziom} B · angiologia / stany pilne · \textbf{cechy:} pilne-skierowanie, bez-lekow, lateralizacja, wartosci-liczbowe, odstawienie\par
\textit{Scenariusz.} 58-latka na HTZ, po długim locie, ból i obrzęk lewej łydki (39 vs 36 cm), Wells 3. Bez leków, pilne skierowanie na SOR (USG doppler, D-dimery), wstrzymanie HTZ, kontrola po wypisie.\par
\textit{Ocena recenzenta.} Oba projekty bez zmian: pusta recepta, pilne skierowanie na SOR „dziś” na pierwszym miejscu, wstrzymanie HTZ i telefon 112 przy objawach zatorowości. Werdykt: \textbf{gotowy bez zmian}; oceny 1--5: dokładność 5, kompletność 5, zwięzłość 5, zrozumiałość dla pacjenta 5, spójność wewnętrzna 5.\par
\textit{Mocne strony:}
\begin{itemize}[nosep]
\item Pilność skierowania oddana wprost („dziś, bez zwlekania”), kontrola po wypisie z SOR oddzielona od skierowania.
\item Obwody łydek (lewa 39 cm, prawa 36 cm) bez zamiany; lateralizacja spójna.
\item Wstrzymanie hormonalnej terapii zastępczej w zaleceniach.
\end{itemize}
\textit{Poprawki lekarza:} brak – oba projekty zatwierdzone bez zmian.\par

\subsection*{111. Zaburzenia lękowe uogólnione z bezsennością}\phantomsection\addcontentsline{toc}{subsection}{111. Zaburzenia lękowe uogólnione z bezsennością}\label{wizyta:111}
\textbf{Poziom} B · psychiatria · \textbf{cechy:} zaprzeczenia, brak-ilosci, wartosci-liczbowe\par
\textit{Scenariusz.} 31-latek z 3-miesięcznym lękiem i bezsennością, GAD-7 12 pkt, bez myśli samobójczych. Hydroksyzyna doraźnie (bez podanej ilości), higiena snu, skierowanie na CBT-I, numery kryzysowe. Dyktowanie bez nagłówka „Recepta”.\par
\textit{Ocena recenzenta.} Projekty kompletne, z poprawnym numerem kryzysowym i zachowanym zaprzeczeniem myśli samobójczych, ale zalecenie o prowadzeniu samochodu znów zmieniło sens („dopóki lekarz nie oceni” zamiast „dopóki pan nie zobaczy”). Werdykt: \textbf{istotne poprawki}; oceny 1--5: dokładność 3, kompletność 5, zwięzłość 4, zrozumiałość dla pacjenta 4, spójność wewnętrzna 5.\par
\textit{Mocne strony:}
\begin{itemize}[nosep]
\item Numer Centrum Wsparcia złożony poprawnie z dyktowania słownego (800 70 2222), razem ze 112.
\item Zaprzeczenie myśli rezygnacyjnych i samobójczych zachowane; GAD-7 odnotowany.
\item Brak ilości hydroksyzyny oznaczony przez system.
\end{itemize}
\textit{Braki:}
\begin{itemize}[nosep]
\item Zalecenie dotyczące prowadzenia samochodu przypisane ocenie lekarza zamiast obserwacji pacjenta.
\item Dopisana jednostka „filiżanek”.
\end{itemize}
\textit{Poprawki lekarza:}
\begin{itemize}[nosep]
\item etap 1 (recepta i zalecenia), \texttt{prescription.0.quantity}: recepta / umiarkowany / źródło: dyktowanie / wykrył: uwaga systemu. Nie podałem ilości hydroksyzyny (system to oznaczył); uzupełniam 1 opakowanie 30 tabletek.
\item etap 1 (recepta i zalecenia), \texttt{patient\_instructions.recommendations.1}: błędny fakt / umiarkowany / źródło: system / wykrył: lekarz. Podyktowałem „dopóki pan nie zobaczy, jak lek działa”; projekt zmienił to na „dopóki lekarz nie oceni” – inne zalecenie co do prowadzenia samochodu. Przywracam sens dyktowania.
\item etap 2 (karta wizyty), \texttt{visit\_note.history}: treść spoza dyktowania / kosmetyczny / źródło: system / wykrył: lekarz. Dopisana jednostka „filiżanek” przy kawie, której nie było w dyktowaniu.
\end{itemize}

\subsection*{112. Choroba zwyrodnieniowa prawego stawu kolanowego}\phantomsection\addcontentsline{toc}{subsection}{112. Choroba zwyrodnieniowa prawego stawu kolanowego}\label{wizyta:112}
\textbf{Poziom} B · ortopedia / reumatologia · \textbf{cechy:} postac-nietabletkowa, jednostki-slownie, lateralizacja, zaprzeczenia\par
\textit{Scenariusz.} 68-latek z chorobą wrzodową w wywiadzie i gonartrozą prawego kolana. Ketoprofen w żelu 2,5\% zamiast NLPZ doustnych, paracetamol, RTG prawego kolana, fizjoterapia.\par
\textit{Ocena recenzenta.} Oba projekty bez zmian: żel z ketoprofenem podany w procentach dopasowany do 25 mg/g, strona prawa spójna we wszystkich dokumentach, ostrzeżenie przed doustnymi NLPZ zachowane. Werdykt: \textbf{gotowy bez zmian}; oceny 1--5: dokładność 5, kompletność 5, zwięzłość 5, zrozumiałość dla pacjenta 5, spójność wewnętrzna 5.\par
\textit{Mocne strony:}
\begin{itemize}[nosep]
\item „Dwa i pół procent” dopasowane do żelu 25 mg/g.
\item Prawe kolano w dawkowaniu, skierowaniu na RTG, wywiadzie, badaniu i rozpoznaniu; lewe opisane jako prawidłowe.
\item Żaden doustny NLPZ na recepcie; ostrzeżenie (choroba wrzodowa) w zaleceniach.
\end{itemize}
\textit{Poprawki lekarza:} brak – oba projekty zatwierdzone bez zmian.\par

\subsection*{113. Nudności i wymioty ciężarnych w 10. tygodniu ciąży}\phantomsection\addcontentsline{toc}{subsection}{113. Nudności i wymioty ciężarnych w 10. tygodniu ciąży}\label{wizyta:113}
\textbf{Poziom} B · położnictwo · \textbf{cechy:} ciaza, nazwa-handlowa, lek-zlozony, kontynuacja-bez-recepty\par
\textit{Scenariusz.} 29-latka w 10. tygodniu pierwszej ciąży, PUQE 9 pkt, ślad ciał ketonowych. Xonvea ze schematem warunkowym, kwas foliowy kontynuowany (bez recepty).\par
\textit{Ocena recenzenta.} Zalecenia i karta bez zarzutu, warunkowe dawkowanie zachowane; jedyna poprawka to ręczny wybór Xonvei po fałszywym konflikcie nazwy handlowej (substancje połączone spójnikiem „z”). Werdykt: \textbf{istotne poprawki}; oceny 1--5: dokładność 4, kompletność 5, zwięzłość 5, zrozumiałość dla pacjenta 5, spójność wewnętrzna 5.\par
\textit{Mocne strony:}
\begin{itemize}[nosep]
\item Warunkowe dawkowanie (2 tabletki wieczorem, ewentualnie dodatkowo 1 rano) przepisane dokładnie.
\item Kwas foliowy („kontynuować”) poza receptą, w zaleceniach.
\item Objawy alarmowe ze wskazaniem szpitala; PUQE i ciała ketonowe w karcie.
\end{itemize}
\textit{Braki:}
\begin{itemize}[nosep]
\item Fałszywy konflikt nazwy handlowej dla Xonvei – ten sam wzorzec co przy Symbicorcie (108).
\end{itemize}
\textit{Poprawki lekarza:}
\begin{itemize}[nosep]
\item etap 1 (recepta i zalecenia), \texttt{prescription.0.product}: recepta / umiarkowany / źródło: system / wykrył: uwaga systemu. Fałszywy konflikt nazwy handlowej: substancje zapisane „doksylamina z pirydoksyną” (spójnik „z”), a Xonvea to właśnie doksylamina + pirydoksyna. Wybieram Xonvea.
\end{itemize}

\subsection*{114. Ostry nieżyt żołądkowo-jelitowy z łagodnym odwodnieniem u 2,5-latka}\phantomsection\addcontentsline{toc}{subsection}{114. Ostry nieżyt żołądkowo-jelitowy z łagodnym odwodnieniem u 2,5-latka}\label{wizyta:114}
\textbf{Poziom} B · pediatria · \textbf{cechy:} pediatria, dawka-wagowa, postac-nietabletkowa, wartosci-liczbowe\par
\textit{Scenariusz.} Chłopiec 12 kg (tydzień temu 12,5 kg), biegunka i wymioty, odwodnienie ok. 4\%. Doustny płyn nawadniający (Gastrolit, w rejestrze bez mocy), paracetamol w czopkach 125 mg.\par
\textit{Ocena recenzenta.} Projekty wierne i kompletne; jedyna interwencja to ręczny wybór Gastrolitu, który w rejestrze nie ma mocy (system słusznie poprosił o wybór). Werdykt: \textbf{istotne poprawki}; oceny 1--5: dokładność 5, kompletność 5, zwięzłość 5, zrozumiałość dla pacjenta 5, spójność wewnętrzna 5.\par
\textit{Mocne strony:}
\begin{itemize}[nosep]
\item Paracetamol dopasowany do czopków 125 mg, nie do innej postaci.
\item Dwie masy ciała (12 kg obecnie, 12,5 kg tydzień temu) i dwie temperatury (38,9 w wywiadzie, 38,2 w badaniu) przypisane poprawnie.
\item Objawy odwodnienia wymagające szpitala wymienione kompletnie, ze wskazaniem szpitala.
\end{itemize}
\textit{Braki:}
\begin{itemize}[nosep]
\item Produkt bez mocy w rejestrze nie może być wybrany automatycznie nawet przy jednoznacznej nazwie handlowej.
\end{itemize}
\textit{Poprawki lekarza:}
\begin{itemize}[nosep]
\item etap 1 (recepta i zalecenia), \texttt{prescription.0.product}: recepta / umiarkowany / źródło: rejestr / wykrył: uwaga systemu. Gastrolit ma w rejestrze moc „-” (produkt złożony), więc system bez mocy nie wybiera produktu; jedyny Gastrolit na liście – wybieram go.
\end{itemize}

\subsection*{115. Ostre bakteryjne zapalenie zatok – dyktowanie z rozpoznawania mowy}\phantomsection\addcontentsline{toc}{subsection}{115. Ostre bakteryjne zapalenie zatok – dyktowanie z rozpoznawania mowy}\label{wizyta:115}
\textbf{Poziom} C · laryngologia · \textbf{cechy:} dyktowanie-chaotyczne, blad-rozpoznawania-mowy, lek-zlozony, autokorekta, nazwa-handlowa, jednostki-slownie, postac-nietabletkowa\par
\textit{Scenariusz.} 45-latek, przebieg dwufazowy z gorączką i ropną wydzieliną – kryteria bakteryjnego zapalenia zatok. Dyktowanie jak z transkrypcji mowy: wtrącenia, autokorekty (7 $\rightarrow$ 10 dni, 1 $\rightarrow$ 2 opakowania), nazwa handlowa zapisana fonetycznie („amoksyklaw”). Amoksycylina z kwasem klawulanowym i mometazon donosowo.\par
\textit{Ocena recenzenta.} Mimo chaotycznego dyktowania (wtrącenia, autokorekty, fonetycznie zapisana nazwa handlowa) recepta i zalecenia bez zmian; w karcie powtórzony błąd rodzaju gramatycznego („zatokiem”). Werdykt: \textbf{drobne poprawki}; oceny 1--5: dokładność 5, kompletność 5, zwięzłość 5, zrozumiałość dla pacjenta 5, spójność wewnętrzna 5.\par
\textit{Mocne strony:}
\begin{itemize}[nosep]
\item „Amoksyklaw” rozpoznany jako Amoksiklav 875 + 125 mg.
\item Obie autokorekty uwzględnione: 10 dni zamiast 7 i dwa opakowania zamiast jednego.
\item Wtrącenia nie trafiły do dokumentów; „nie wiem jakie” oddane neutralnie jako „nazwy nieznane”.
\end{itemize}
\textit{Braki:}
\begin{itemize}[nosep]
\item Błąd rodzaju gramatycznego w badaniu.
\end{itemize}
\textit{Poprawki lekarza:}
\begin{itemize}[nosep]
\item etap 2 (karta wizyty), \texttt{visit\_note.physical\_examination}: forma/język / kosmetyczny / źródło: system / wykrył: lekarz. Błąd językowy: „nad prawym zatokiem szczękowym; lewy zatok” zamiast „nad prawą zatoką szczękową; lewa zatoka”.
\end{itemize}

\subsection*{116. Angina paciorkowcowa – amoksycylina u pacjentki uczulonej na amoksycylinę}\phantomsection\addcontentsline{toc}{subsection}{116. Angina paciorkowcowa – amoksycylina u pacjentki uczulonej na amoksycylinę}\label{wizyta:116}
\textbf{Poziom} C · choroby zakaźne / bezpieczeństwo · \textbf{cechy:} konflikt-alergii, dyktowanie-chaotyczne, brak-ilosci, autokorekta, wartosci-liczbowe\par
\textit{Scenariusz.} 36-latka z anginą paciorkowcową; w karcie alergia na amoksycylinę (pokrzywka, obrzęk warg). Lekarz w pośpiechu dyktuje amoksycylinę i deklaruje „nic jej nie uczula”. Sprawdza, czy system (który niczego nie decyduje) w ogóle sygnalizuje sprzeczność z kartą pacjenta.\par
\textit{Ocena recenzenta.} Projekty wiernie oddały dyktowanie, łącznie z błędem lekarza: amoksycylina dla pacjentki uczulonej na amoksycylinę, bez żadnego ostrzeżenia systemu. Karta tym razem jawnie odnotowała sprzeczność między wypowiedzią pacjentki a kartą. Werdykt: \textbf{istotne poprawki}; oceny 1--5: dokładność 4, kompletność 5, zwięzłość 5, zrozumiałość dla pacjenta 5, spójność wewnętrzna 3.\par
\textit{Mocne strony:}
\begin{itemize}[nosep]
\item Recepta i zalecenia zgodne z dyktowaniem mimo wtrąceń; brak ilości amoksycyliny oznaczony.
\item Wywiad zestawia wprost zaprzeczenie pacjentki i alergię z karty, bez cichego „godzenia” sprzeczności.
\item Autokorekta „nie pytałem, nie, przepraszam…” oddana poprawnie.
\end{itemize}
\textit{Braki:}
\begin{itemize}[nosep]
\item Brak jakiejkolwiek uwagi systemu o sprzeczności przepisywanego leku z alergią w karcie – błąd istotny wykrył dopiero lekarz.
\item Uwaga weryfikatora błędnie odczytała autokorektę lekarza (fałszywy alarm).
\end{itemize}
\textit{Poprawki lekarza:}
\begin{itemize}[nosep]
\item etap 1 (recepta i zalecenia), \texttt{prescription.0}: recepta / istotny / źródło: dyktowanie / wykrył: lekarz. W karcie pacjentki alergia na amoksycylinę (pokrzywka, obrzęk warg – reakcja natychmiastowa); podyktowałem amoksycylinę przez pomyłkę. System nie zasygnalizował sprzeczności z kartą. Zmieniam na klarytromycynę 250 mg 2 $\times$ dziennie przez 10 dni.
\item etap 1 (recepta i zalecenia), \texttt{prescription.0.quantity}: recepta / umiarkowany / źródło: dyktowanie / wykrył: uwaga systemu. Nie podałem ilości antybiotyku (system to oznaczył); uzupełniam 20 tabletek.
\end{itemize}

\subsection*{117. Zaostrzenie niewydolności serca u 79-latka – wiele zmian leczenia}\phantomsection\addcontentsline{toc}{subsection}{117. Zaostrzenie niewydolności serca u 79-latka – wiele zmian leczenia}\label{wizyta:117}
\textbf{Poziom} C · kardiologia / geriatria · \textbf{cechy:} polipragmazja, zmiana-dawki, odstawienie, kontynuacja-bez-recepty, przedluzenie-recepty, wiek-podeszly, wartosci-liczbowe\par
\textit{Scenariusz.} 79-latek z HFrEF (EF 35\%), przyrost masy 3 kg, NYHA III, prawdopodobnie po NLPZ. W jednym dyktowaniu: zwiększenie furosemidu na tydzień, nowy spironolakton, bisoprolol bez zmian (bez recepty), przedłużenie ramiprilu, odstawienie ibuprofenu, kontrola potasu.\par
\textit{Ocena recenzenta.} Recepta przy czterech jednoczesnych zmianach leczenia bezbłędna; słabe były zalecenia (brak podyktowanych zmian leczenia, w tym kontynuacji bisoprololu) i karta (dawka ibuprofenu 402 mg, pominięty plan leczenia). Werdykt: \textbf{istotne poprawki}; oceny 1--5: dokładność 3, kompletność 3, zwięzłość 5, zrozumiałość dla pacjenta 3, spójność wewnętrzna 4.\par
\textit{Mocne strony:}
\begin{itemize}[nosep]
\item Na recepcie dokładnie trzy właściwe leki: czasowo zwiększony furosemid, nowy spironolakton, przedłużony ramipril; bisoprolol i ibuprofen poza receptą.
\item Postać ramiprilu słusznie oznaczona do wyboru.
\item Karta kompletna co do parametrów (EF 35\%, przyrost 3 kg, NT-proBNP 2400, eGFR 42, NYHA III).
\end{itemize}
\textit{Braki:}
\begin{itemize}[nosep]
\item Zalecenia bez zmian leczenia – u 79-latka bisoprolol „bez zmian” nie pojawia się w żadnym dokumencie dla pacjenta.
\item „Czterysta, dwa, trzy razy dziennie” odczytane jako dawka 402 mg; weryfikator tego nie zgłosił.
\item Pominięty plan „leczenie ambulatoryjne z intensyfikacją diuretyku” (wskazał weryfikator).
\end{itemize}
\textit{Poprawki lekarza:}
\begin{itemize}[nosep]
\item etap 1 (recepta i zalecenia), \texttt{prescription.2.product}: recepta / umiarkowany / źródło: rejestr / wykrył: uwaga systemu. Ramipril 5 mg występuje jako tabletki i kapsułki, postaci nie podałem – system słusznie nie wybrał; wybieram tabletki.
\item etap 1 (recepta i zalecenia), \texttt{patient\_instructions.recommendations}: pominięcie / umiarkowany / źródło: system / wykrył: lekarz. Zalecenia pomijają podyktowane zmiany leczenia: czasowe zwiększenie furosemidu, nowy spironolakton i „bisoprolol bez zmian” (bisoprololu nie ma na recepcie, więc pacjent mógłby go odstawić). Dopisuję je na początku zaleceń.
\item etap 2 (karta wizyty), \texttt{visit\_note.history}: błędny fakt / umiarkowany / źródło: system / wykrył: lekarz. „Ibuprofen czterysta, dwa, trzy razy dziennie” to 400 mg 2–3 razy dziennie; model wpisał dawkę 402 mg. Weryfikator tego nie zgłosił. Przy okazji: „wcześniej wchodził na trzecie” zmienione na „wcześniej duszność na III piętrze” – przywracam sens dyktowania.
\item etap 2 (karta wizyty), \texttt{visit\_note.diagnosis}: pominięcie / umiarkowany / źródło: system / wykrył: weryfikator. Pominięty podyktowany plan: leczenie ambulatoryjne z intensyfikacją diuretyku (wskazał weryfikator).
\end{itemize}

\subsection*{118. Pozaszpitalne zapalenie płuc u 2,5-letniej dziewczynki – dawkowanie wagowe}\phantomsection\addcontentsline{toc}{subsection}{118. Pozaszpitalne zapalenie płuc u 2,5-letniej dziewczynki – dawkowanie wagowe}\label{wizyta:118}
\textbf{Poziom} C · pediatria · \textbf{cechy:} pediatria, dawka-wagowa, wartosci-liczbowe, postac-nietabletkowa, autokorekta, brak-ilosci, dyktowanie-chaotyczne, lateralizacja\par
\textit{Scenariusz.} Dziewczynka 12 kg, lewostronne zapalenie płuc. Amoksycylina 90 mg/kg/d: lekarz liczy na głos (1080 mg/d $\rightarrow$ 540 mg co 12 h $\rightarrow$ 10,8 ml) i ostatecznie podaje 10,5 ml. Ibuprofen w zawiesinie bez ilości.\par
\textit{Ocena recenzenta.} Bardzo dobre projekty przy trudnym dyktowaniu: z liczenia dawki na głos zachowana ostateczna objętość 10,5 ml, zawiesiny dopasowane poprawnie, strona lewa bez pomyłki; jedyna poprawka to brakująca ilość ibuprofenu z dyktowania. Werdykt: \textbf{istotne poprawki}; oceny 1--5: dokładność 5, kompletność 5, zwięzłość 5, zrozumiałość dla pacjenta 5, spójność wewnętrzna 5.\par
\textit{Mocne strony:}
\begin{itemize}[nosep]
\item Ostateczna objętość amoksycyliny 10,5 ml (nie 10,8 ml z obliczeń); zawiesina 250 mg/5 ml.
\item Brak ilości ibuprofenu oznaczony przez system.
\item Karta bez nadinterpretacji („pije”), strona lewa, dwie temperatury i parametry oddechowe przepisane bez błędów.
\end{itemize}
\textit{Poprawki lekarza:}
\begin{itemize}[nosep]
\item etap 1 (recepta i zalecenia), \texttt{prescription.1.quantity}: recepta / umiarkowany / źródło: dyktowanie / wykrył: uwaga systemu. Nie podałem ilości ibuprofenu (system to oznaczył); uzupełniam 1 butelkę 100 ml.
\end{itemize}

\subsection*{119. Skręcenie prawego stawu skokowego – autokorekta strony i leczenia}\phantomsection\addcontentsline{toc}{subsection}{119. Skręcenie prawego stawu skokowego – autokorekta strony i leczenia}\label{wizyta:119}
\textbf{Poziom} C · ortopedia / urazy · \textbf{cechy:} lateralizacja, autokorekta, zaprzeczenia, dyktowanie-chaotyczne, postac-nietabletkowa, nazwa-handlowa\par
\textit{Scenariusz.} 26-latek po urazie inwersyjnym, kryteria ottawskie ujemne (bez RTG). Lekarz w etapie 1 odwołuje ibuprofen i zwolnienie lekarskie, w etapie 2 poprawia stronę („lewą… nie, prawą”). Żel z diklofenakiem podany nazwą handlową bez mocy.\par
\textit{Ocena recenzenta.} Autokorekty w obu etapach rozwiązane poprawnie (strona prawa, rezygnacja z ibuprofenu i zwolnienia), karta kompletna łącznie z decyzją „bez wskazań do RTG”; jedyna poprawka to wybór produktu przy niepodanej mocy żelu. Werdykt: \textbf{istotne poprawki}; oceny 1--5: dokładność 5, kompletność 5, zwięzłość 5, zrozumiałość dla pacjenta 4, spójność wewnętrzna 5.\par
\textit{Mocne strony:}
\begin{itemize}[nosep]
\item Korekta strony „lewą… nie, prawą” zastosowana konsekwentnie; lewy staw opisany jako bez zmian.
\item Odwołany ibuprofen i odwołane zwolnienie nie pojawiły się w żadnym dokumencie.
\item Kryteria ottawskie z decyzją „brak wskazań do RTG”; brak mocy żelu oznaczony do wyboru.
\end{itemize}
\textit{Braki:}
\begin{itemize}[nosep]
\item Objawy alarmowe bez słowa „pilnie” (znaczenie niesie samo pole, ale pacjent nie dostaje polecenia działania).
\end{itemize}
\textit{Poprawki lekarza:}
\begin{itemize}[nosep]
\item etap 1 (recepta i zalecenia), \texttt{prescription.0.product}: recepta / umiarkowany / źródło: dyktowanie / wykrył: uwaga systemu. Nie podałem mocy żelu; system słusznie poprosił o wybór. Jedyny kandydat to Voltaren MAX 23,2 mg/g – wybieram go.
\end{itemize}

\subsection*{120. Półpasiec oczny z objawem Hutchinsona – pilne skierowanie do okulisty}\phantomsection\addcontentsline{toc}{subsection}{120. Półpasiec oczny z objawem Hutchinsona – pilne skierowanie do okulisty}\label{wizyta:120}
\textbf{Poziom} C · dermatologia / okulistyka / stany pilne · \textbf{cechy:} pilne-skierowanie, wiek-podeszly, wartosci-liczbowe, zaprzeczenia, dyktowanie-chaotyczne, lateralizacja\par
\textit{Scenariusz.} 79-latka z półpaścem w zakresie V1 po prawej, objaw Hutchinsona dodatni. Walacyklowir, paracetamol, pilne skierowanie do okulisty tego samego dnia, kontrola za tydzień po konsultacji. Dyktowanie z wtrąceniami.\par
\textit{Ocena recenzenta.} Najważniejszy element wizyty – pilne skierowanie na ostry dyżur okulistyczny tego samego dnia – zachowany wprost, recepta dopasowana automatycznie, karta bez zarzutu; w zaleceniach pominięte „natychmiast do szpitala” przy objawach ocznych i „zaczynamy dziś” przy walacyklowirze. Werdykt: \textbf{istotne poprawki}; oceny 1--5: dokładność 5, kompletność 4, zwięzłość 5, zrozumiałość dla pacjenta 4, spójność wewnętrzna 5.\par
\textit{Mocne strony:}
\begin{itemize}[nosep]
\item Pilne skierowanie „dziś” na ostry dyżur okulistyczny na pierwszym miejscu zaleceń; kontrola po okuliście oddzielona.
\item Moce z jednostkami (1000 mg, 500 mg) – produkty wybrane automatycznie.
\item Karta kompletna mimo wtrąceń: objaw Hutchinsona, zakres V1, strona prawa, zaprzeczenia, eGFR 62.
\end{itemize}
\textit{Braki:}
\begin{itemize}[nosep]
\item Pominięte miejsce zgłoszenia przy objawach alarmowych (szpital) – powtarzalny błąd (108, 109).
\item Pominięte „zaczynamy dziś”.
\item Weryfikator zgłosił dopisanie strony „prawe” przy łzawieniu – wniosek zgodny z kontekstem, fałszywy alarm.
\end{itemize}
\textit{Poprawki lekarza:}
\begin{itemize}[nosep]
\item etap 1 (recepta i zalecenia), \texttt{prescription.0.dosage}: pominięcie / kosmetyczny / źródło: system / wykrył: lekarz. Pominięte „zaczynamy dziś” – przy półpaścu ważne rozpoczęcie leczenia od razu.
\item etap 1 (recepta i zalecenia), \texttt{patient\_instructions.warning\_signs}: pominięcie / umiarkowany / źródło: system / wykrył: lekarz. Pominięte „natychmiast do szpitala” przy objawach zajęcia oka – pacjentka nie wie, gdzie i jak pilnie się zgłosić.
\end{itemize}


\par}

\section{Dane i~odtworzenie}
{\sloppy
\begin{itemize}
  \item \texttt{data/visits/visit\_101.json} \ldots{} \texttt{visit\_120.json} -- scenariusze (dyktowania, poziom, cechy, hipotezy);
        \texttt{data/patients.json} -- dane pacjentów.
  \item \texttt{results/v4\_doctor/} -- wyniki przebiegu: \texttt{visits.csv} (wizyta), \texttt{corrections.csv} (poprawka),
        \texttt{field\_changes.csv} (zmienione pole), \texttt{issues.csv} (uwaga systemu), \texttt{features.csv} (wizyta $\times$ cecha),
        \texttt{hypotheses.csv}, \texttt{run.json}; w~\texttt{visits/<wizyta>/} projekty i~wersje zatwierdzone obu etapów,
        \texttt{log.json} i~\texttt{analysis.json}.
\end{itemize}
\par}
Przebieg (lekarz zatwierdza etapy plikami \texttt{stageN.done} albo w~terminalu), zestawienie z~eksportem i~fragmenty tego raportu:
\begin{verbatim}
python benchmark.py
python benchmark.py --summary-only --export results/v4_doctor
python benchmark_report.py results/v4_doctor docs/benchmark/generated
\end{verbatim}

\end{document}
